\part{Durchlaufen}

\chapter{Eulerzüge}

\section{Zurück nach Königsberg}

Die Frage der Bürger lautete: Gibt es einen Spaziergang, der jede der sieben Brücken genau einmal überquert?

Aus Kapitel~1 wissen wir schon, wie der zugehörige Graph aussieht --- und dass er kein Graph im Sinne von Definition 1.1 ist, weil zwischen Insel und Nordufer zwei Brücken verlaufen. Bevor wir die Frage beantworten können, müssen wir den Begriff also erweitern. Das war in Kapitel~1 angekündigt, und hier ist die Stelle.

\begin{definition}
\boxtitelm{Definition 6.1 --- Multigraph}
Ein \Begriff{Multigraph} ist ein Tripel $G = (V, E, \varphi)$ aus
\begin{itemize}[leftmargin=1.4em, itemsep=3pt]
\item einer endlichen, nichtleeren Menge $V$ von Ecken,
\item einer endlichen Menge $E$ von Kanten,
\item einer Abbildung $\varphi$, die jeder Kante $e \in E$ ein Paar $\{u,v\}$ von Ecken zuordnet.
\end{itemize}
Ist $\varphi(e) = \{u,u\}$, so heißt $e$ eine \Begriff{Schlinge}. Zwei verschiedene Kanten mit demselben Bild heißen \Begriff{parallel}.

Der \Begriff{Grad} $\deg(v)$ ist die Anzahl der Kantenenden an $v$; eine Schlinge an $v$ zählt dabei zweimal.
\end{definition}

\begin{merksatz}
\boxtitel{Warum die Kanten jetzt einen eigenen Namen brauchen}
Bei Definition 1.1 \kursiv{war} eine Kante die Menge ihrer beiden Endecken. Zwei Kanten mit denselben Enden wären dann dieselbe Kante --- Parallelen sind unmöglich.

Deshalb trennt Definition 6.1 zwei Dinge, die vorher zusammenfielen: die Kante als eigenständiges Ding (ein Element von $E$) und ihre Lage (das Bild unter $\varphi$). Zwei Brücken zwischen Insel und Nordufer sind zwei verschiedene Elemente von $E$ mit demselben Bild.

\begin{center}
\begin{minipage}{0.88\textwidth}
\kursiv{Das ist ein Muster, das in der Mathematik oft vorkommt: Ein Begriff wird erweitert, indem man eine Identifizierung aufgibt. Alles Bisherige bleibt gültig --- ein Graph nach Definition 1.1 ist ein Multigraph ohne Schlingen und ohne Parallelen.}
\end{minipage}
\end{center}
\end{merksatz}

\begin{merksatz}
\boxtitel{Das Handschlaglemma gilt weiter}
Satz 1.3 bleibt in der Fassung
\[ \sum_{v \in V} \deg(v) = 2\,|E| \]
richtig, und der Beweis ändert sich kaum: Jede Kante hat zwei Enden, jedes Ende liegt an genau einer Ecke.

Die Schlingenregel \enquote{zählt doppelt} ist genau das, was diesen Beweis funktionieren lässt: Eine Schlinge hat zwei Enden, und beide liegen an derselben Ecke. Hätte man Schlingen einfach gezählt, wäre die Formel falsch.
\end{merksatz}

\begin{merksatz}
\boxtitel{Der Königsberger Multigraph}
\begin{center}
\begin{tikzpicture}[scale=1.4]
\node[knotenb] (N) at (0,1.1) {$N$};
\node[knotenb] (S) at (0,-1.1) {$S$};
\node[knotenb] (I) at (-1.2,0) {$I$};
\node[knotenb] (O) at (1.2,0) {$O$};
\draw[kante] (I) to[bend left=34] (N);
\draw[kante] (I) to[bend right=34] (N);
\draw[kante] (I) to[bend left=34] (S);
\draw[kante] (I) to[bend right=34] (S);
\draw[kante] (I) -- (O);
\draw[kante] (N) -- (O);
\draw[kante] (S) -- (O);
\end{tikzpicture}
\end{center}

\begin{center}
\renewcommand{\arraystretch}{1.45}
\begin{tabular}{@{}P{2.4cm}P{3.6cm}P{1.6cm}@{}}
\toprule
\textbf{Ecke} & \textbf{Kanten} & $\deg$\\
\midrule
Insel $I$ & $2$ nach $N$, $2$ nach $S$, $1$ nach $O$ & $\mathbf{5}$\\
Nordufer $N$ & $2$ nach $I$, $1$ nach $O$ & $\mathbf{3}$\\
Südufer $S$ & $2$ nach $I$, $1$ nach $O$ & $\mathbf{3}$\\
Ostland $O$ & je $1$ nach $I$, $N$, $S$ & $\mathbf{3}$\\
\bottomrule
\end{tabular}
\end{center}

Probe: $5+3+3+3 = 14 = 2 \cdot 7$. \checkmark
\end{merksatz}

\section{Die Begriffe}

\begin{definition}
\boxtitelm{Definition 6.2 --- Eulerzug, Eulerkreis}
Sei $G$ ein Multigraph.
\begin{itemize}[leftmargin=1.4em, itemsep=4pt]
\item Ein \Begriff{Eulerzug} ist ein Kantenzug, der \kursiv{jede Kante von $G$ genau einmal} benutzt.
\item Ein \Begriff{Eulerkreis} ist ein Eulerzug, der an derselben Ecke endet, an der er begann.
\item $G$ heißt \Begriff{eulersch}, wenn er einen Eulerkreis besitzt.
\end{itemize}
\end{definition}

\begin{warnung}
\boxtitelr{Was hier \kursiv{nicht} gefordert wird}
Ein Eulerzug darf Ecken beliebig oft besuchen. Nur die \kursiv{Kanten} dürfen sich nicht wiederholen.

Das ist der Unterschied zum Weg (Definition 2.2), bei dem die Ecken verschieden sein müssen. Und es ist der Unterschied zum Hamiltonkreis des nächsten Kapitels, bei dem umgekehrt jede Ecke genau einmal besucht wird und Kanten übrig bleiben dürfen.

\begin{center}
\renewcommand{\arraystretch}{1.5}
\begin{tabular}{@{}P{2.6cm}P{2.6cm}P{2.6cm}@{}}
\toprule
& \textbf{Jede Kante} & \textbf{Jede Ecke}\\
\midrule
Eulerkreis & \textcolor{gruen}{genau einmal} & beliebig oft\\
Hamiltonkreis & darf fehlen & \textcolor{gruen}{genau einmal}\\
\bottomrule
\end{tabular}
\end{center}

\begin{center}
\begin{minipage}{0.88\textwidth}
\kursiv{Die beiden Begriffe sehen einander zum Verwechseln ähnlich. Sie sind es nicht. Für den einen gibt es ein einfaches, vollständiges Kriterium --- Satz 6.4. Für den anderen kennt niemand eines, und es gibt gute Gründe zu glauben, dass es keines gibt.}
\end{minipage}
\end{center}
\end{warnung}

\section{Die notwendige Bedingung}

\begin{merksatz}
\boxtitel{Satz 6.3 --- Was ein Eulerkreis erzwingt}
Besitzt ein Multigraph $G$ einen Eulerkreis und hat er mindestens eine Kante, so gilt:
\begin{itemize}[leftmargin=1.4em, itemsep=2pt]
\item alle Ecken mit Grad $\ge 1$ liegen in einer einzigen Zusammenhangskomponente,
\item jede Ecke hat geraden Grad.
\end{itemize}
\end{merksatz}

\begin{merksatz}
\boxtitel{Beweis}
\schritt{Erster Teil --- der Zusammenhang}
Der Eulerkreis benutzt jede Kante. Also berührt er jede Ecke, an der eine Kante hängt. Ein Kantenzug verläuft aber innerhalb einer Komponente --- er kann nie von einer Komponente in eine andere springen, weil jeder Schritt eine Kante benutzt. Also liegen alle Ecken mit Grad $\ge 1$ in derselben Komponente.

\schritt{Zweiter Teil --- die Parität}
Sei $v$ eine beliebige Ecke und $Z$ der Eulerkreis. Wir zählen die Kantenenden an $v$ auf zwei Weisen.

\kursiv{Erste Zählung.} Es sind $\deg(v)$ Stück --- das ist die Definition.

\kursiv{Zweite Zählung.} Wir verfolgen $Z$ und zählen, wie oft er die Ecke $v$ berührt.

Jedes Mal, wenn $Z$ die Ecke $v$ als \kursiv{Zwischenecke} durchläuft, kommt er über eine Kante an und geht über eine andere weg. Das verbraucht zwei Kantenenden an $v$.

Bleibt der Start. Der Eulerkreis beginnt an einer Ecke $r$ und endet dort. Ist $v \neq r$, so gibt es keinen weiteren Beitrag: Alle Berührungen sind Durchläufe, also ist $\deg(v)$ das Doppelte der Anzahl der Durchläufe --- gerade.

Ist $v = r$, so kommt zusätzlich das allererste Kantenende (beim Aufbruch) und das allerletzte (bei der Rückkehr) hinzu. Das sind zwei weitere Enden. Also ist auch $\deg(r)$ gerade.

\schritt{Warum jedes Kantenende genau einmal gezählt wird}
Jede Kante an $v$ wird von $Z$ genau einmal benutzt --- das ist die Eulereigenschaft. Bei dieser Benutzung wird genau ein Ende an $v$ verbraucht (oder bei einer Schlinge beide). Es bleibt also kein Ende ungezählt und keines wird doppelt gezählt. $\blacksquare$

\begin{center}
\begin{tikzpicture}[scale=1.2]
\node[knotenb] (v) at (0,0) {$v$};
\draw[kanted, -{Stealth[length=5pt]}] (-1.5,0.75) -- (v);
\draw[kanted, -{Stealth[length=5pt]}] (v) -- (1.5,0.75);
\draw[kante, -{Stealth[length=5pt]}] (-1.5,-0.75) -- (v);
\draw[kante, -{Stealth[length=5pt]}] (v) -- (1.5,-0.75);
\node[font=\scriptsize\itshape, color=grau] at (0,-1.35) {jeder Durchlauf verbraucht zwei Enden};
\end{tikzpicture}
\end{center}
\end{merksatz}

\begin{merksatz}
\boxtitel{Die Antwort für Königsberg}
Der Königsberger Multigraph hat vier Ecken mit den Graden $5, 3, 3, 3$ --- \kursiv{alle vier ungerade}.

Nach Satz 6.3 kann er keinen Eulerkreis besitzen. Und auch keinen offenen Eulerzug, wie Satz 6.5 gleich zeigen wird: Dafür dürfte es höchstens zwei ungerade Ecken geben.

\begin{center}
\begin{minipage}{0.9\textwidth}
\textbf{Der gesuchte Spaziergang existiert nicht.} Nicht, weil man ihn noch nicht gefunden hätte --- sondern weil er nicht existieren kann.
\end{minipage}
\end{center}

\begin{center}
\begin{minipage}{0.88\textwidth}
\kursiv{Beachten Sie, worin die Leistung besteht. Die Bürger konnten beliebig lange probieren, ohne je sicher zu sein, ob sie nur den richtigen Weg noch nicht gefunden hatten. Der Beweis beendet das Probieren. Das ist der Unterschied zwischen \enquote{noch keine Lösung gefunden} und \enquote{es gibt keine}.}
\end{minipage}
\end{center}
\end{merksatz}

\section{Die Bedingung genügt auch}

Satz 6.3 sagt, was ein Eulerkreis erzwingt. Der folgende Satz sagt, dass mehr nicht nötig ist --- und das ist die eigentliche Überraschung.

\begin{merksatz}
\boxtitel{Satz 6.4 (Euler 1736, Hierholzer 1873)}
Ein Multigraph $G$ mit mindestens einer Kante besitzt genau dann einen Eulerkreis, wenn
\begin{itemize}[leftmargin=1.4em, itemsep=2pt]
\item alle Ecken mit Grad $\ge 1$ in einer Komponente liegen und
\item jede Ecke geraden Grad hat.
\end{itemize}
\end{merksatz}

\begin{merksatz}
\boxtitel{Beweis der schwierigen Richtung --- Vorbereitung}
Wir dürfen isolierte Ecken weglassen; sie spielen für Kantenzüge keine Rolle. Sei also $G$ zusammenhängend mit lauter geraden Graden.

\schritt{Hilfssatz --- ein Kreis existiert}
Behauptung: Hat jede Ecke von $G$ geraden Grad und gibt es mindestens eine Kante, so enthält $G$ einen Kreis.

Ist eine Schlinge oder ein Paar paralleler Kanten vorhanden, so hat man bereits einen geschlossenen Kantenzug, der keine Kante wiederholt --- das genügt uns. Andernfalls ist $G$ ein gewöhnlicher Graph, in dem jede Ecke mit Grad $\ge 1$ sogar Grad $\ge 2$ hat (gerade und positiv). Nach dem Beweisverfahren aus Satz 4.3 hat ein kreisfreier Graph mit mindestens einer Kante ein Blatt, also eine Ecke vom Grad $1$ --- Widerspruch. Also gibt es einen Kreis.
\end{merksatz}

\begin{merksatz}
\boxtitel{Beweis --- das Verfahren von Hierholzer}
\schritt{Schritt 1 --- einen ersten geschlossenen Zug bauen}
Wir starten an einer beliebigen Ecke $r$ und laufen los, wobei wir keine Kante zweimal benutzen. Wir laufen so lange, bis es nicht mehr weitergeht.

\kursiv{Behauptung: Wir bleiben in $r$ stecken.} Angenommen, wir bleiben in einer Ecke $v \neq r$ stecken. Bis dahin haben wir $v$ einige Male durchlaufen --- jedes Mal zwei Kantenenden verbraucht --- und sind zuletzt angekommen, ohne weitergehen zu können. Das verbraucht eine ungerade Zahl von Enden an $v$. Da $\deg(v)$ gerade ist, bleibt mindestens ein unbenutztes Ende übrig, über das wir weitergehen könnten. Widerspruch.

Also endet der Zug in $r$: Wir haben einen geschlossenen Kantenzug $Z_1$ ohne Kantenwiederholung.

\schritt{Schritt 2 --- Sind schon alle Kanten dabei?}
Wenn ja, ist $Z_1$ ein Eulerkreis, und wir sind fertig.

Wenn nein, sei $H$ der Multigraph, der aus $G$ durch Entfernen aller Kanten von $Z_1$ entsteht (die Ecken bleiben).

\schritt{Schritt 3 --- $H$ hat wieder lauter gerade Grade}
$Z_1$ ist ein geschlossener Kantenzug. Nach demselben Argument wie im Beweis von Satz 6.3 verbraucht er an jeder Ecke eine \kursiv{gerade} Anzahl von Kantenenden.

Gerade minus gerade ist gerade. Also hat in $H$ jede Ecke geraden Grad.

\schritt{Schritt 4 --- $Z_1$ und der Rest berühren sich}
Behauptung: Es gibt eine Ecke $x$, die auf $Z_1$ liegt und in $H$ noch Grad $\ge 1$ hat.

Angenommen nicht. Dann hätte jede Ecke von $Z_1$ in $H$ den Grad $0$. Da $H$ nach Annahme noch Kanten besitzt, gibt es eine Ecke $y$ mit $\deg_H(y) \ge 1$; sie liegt nicht auf $Z_1$.

$G$ ist aber zusammenhängend, es gibt also einen Weg von $r$ nach $y$ in $G$. Auf diesem Weg gibt es eine erste Ecke, die nicht mehr auf $Z_1$ liegt; die Kante, über die man sie erreicht, geht von einer Ecke von $Z_1$ aus und gehört nicht zu $Z_1$ (denn alle Kanten von $Z_1$ verbinden Ecken von $Z_1$, und diese Kante endet außerhalb). Also liegt sie in $H$ --- und ihre Startecke liegt auf $Z_1$ und hat in $H$ Grad $\ge 1$. Widerspruch.

\schritt{Schritt 5 --- Einen zweiten Zug bauen und einsetzen}
Wir wenden Schritt 1 auf $H$ mit Startecke $x$ an und erhalten einen geschlossenen Kantenzug $Z_2$ durch $x$, der nur Kanten aus $H$ benutzt.

Nun setzen wir zusammen: Wir laufen $Z_1$ ab; sobald wir zum ersten Mal bei $x$ ankommen, schieben wir den ganzen Zug $Z_2$ ein und laufen danach $Z_1$ zu Ende.

\begin{center}
\begin{tikzpicture}[scale=1.2]
\draw[kante] (0,0) circle (1.0);
\node[knotenr] (x) at (1.0,0) {$x$};
\draw[kanted] (2.05,0) circle (1.05);
\node[font=\scriptsize\itshape, color=grau] at (-0.9,-1.25) {$Z_1$};
\node[font=\scriptsize\itshape, color=akzent] at (3.0,-1.25) {$Z_2$};
\end{tikzpicture}
\end{center}

Das Ergebnis ist wieder ein geschlossener Kantenzug ohne Kantenwiederholung --- $Z_1$ und $Z_2$ haben ja keine gemeinsame Kante --- und er ist echt länger als $Z_1$.

\schritt{Schritt 6 --- Warum das Verfahren endet}
Bei jeder Wiederholung von Schritt 2 bis 5 wächst die Zahl der benutzten Kanten um mindestens $1$. Da $G$ nur endlich viele Kanten hat, endet das Verfahren nach höchstens $|E|$ Durchgängen. Es endet aber nur im Fall \enquote{alle Kanten dabei} aus Schritt 2 --- also mit einem Eulerkreis. $\blacksquare$
\end{merksatz}

\begin{merksatz}
\boxtitel{Satz 6.5 --- der offene Eulerzug}
Ein zusammenhängender Multigraph besitzt genau dann einen Eulerzug, der nicht geschlossen ist, wenn er \kursiv{genau zwei} Ecken ungeraden Grades hat. Der Zug beginnt dann in der einen und endet in der anderen.

\schritt{Beweis, Richtung 1}
Sei $Z$ ein offener Eulerzug von $a$ nach $b$ mit $a \neq b$. Dieselbe Zählung wie in Satz 6.3 ergibt: Für $v \notin \{a,b\}$ ist $\deg(v)$ gerade. Bei $a$ kommt das Startende hinzu, bei $b$ das Endende --- je eines, also sind $\deg(a)$ und $\deg(b)$ ungerade.

Nach Satz 1.4 ist die Anzahl ungerader Ecken ohnehin gerade; hier sind es genau zwei.

\schritt{Beweis, Richtung 2 --- der Kunstgriff}
Habe $G$ genau zwei ungerade Ecken $a$ und $b$. Wir fügen eine \kursiv{neue} Kante $e^{*}$ zwischen $a$ und $b$ hinzu und nennen den entstandenen Multigraphen $G^{*}$.

In $G^{*}$ sind $\deg(a)$ und $\deg(b)$ um $1$ gewachsen, also gerade; alle übrigen Grade sind unverändert gerade. $G^{*}$ ist zusammenhängend. Nach Satz 6.4 gibt es einen Eulerkreis $Z^{*}$.

Nun lassen wir aus $Z^{*}$ die Kante $e^{*}$ weg. Da ein Eulerkreis geschlossen ist, darf man ihn an jeder Stelle aufschneiden; wir schneiden ihn so auf, dass $e^{*}$ der letzte Schritt ist. Übrig bleibt ein Kantenzug, der alle Kanten von $G$ genau einmal benutzt und von $b$ nach $a$ läuft. $\blacksquare$

\begin{center}
\begin{minipage}{0.88\textwidth}
\kursiv{Der Kunstgriff --- ein Problem durch \kursiv{Hinzufügen} lösen --- ist typisch. Man macht die Aufgabe scheinbar größer und dafür einfacher, löst sie mit einem bereits bewiesenen Satz und nimmt die Zutat am Schluss wieder heraus.}
\end{minipage}
\end{center}
\end{merksatz}

\section{Wie viele Brücken hätten gefehlt?}

\begin{merksatz}
\boxtitel{Königsberg nachgebessert}
Königsberg hat vier ungerade Ecken. Für einen offenen Eulerzug dürften es zwei sein, für einen Eulerkreis null.

\schritt{Eine zusätzliche Brücke genügt für den offenen Zug}
Baut man eine achte Brücke zwischen zwei der ungeraden Ecken --- etwa zwischen Nordufer und Südufer ---, so werden beide gerade. Übrig bleiben die zwei ungeraden Ecken $I$ und $O$, und nach Satz 6.5 gibt es einen Eulerzug von $I$ nach $O$.

\schritt{Zwei zusätzliche Brücken für den Rundgang}
Für einen Eulerkreis müssen alle vier ungeraden Ecken gerade werden. Zwei Brücken genügen: eine zwischen $N$ und $S$, eine zwischen $I$ und $O$. Weniger geht nicht, denn jede Brücke ändert die Parität von höchstens zwei Ecken.

\schritt{Was tatsächlich geschah}
Von den sieben Brücken wurden im Zweiten Weltkrieg fünf zerstört; zwei wurden nicht wieder aufgebaut, und später kam eine Schnellstraßenbrücke hinzu. Der heutige Graph von Kaliningrad hat genau zwei ungerade Ecken --- ein Eulerzug ist möglich, ein Rundgang nicht.
\end{merksatz}

\begin{merksatz}
\boxtitel{Das Problem des Briefträgers}
Muss ein Briefträger jede Straße mindestens einmal begehen und zum Ausgangspunkt zurückkehren, so sucht er einen kürzesten geschlossenen Kantenzug, der jede Kante mindestens einmal enthält.

\begin{center}
\renewcommand{\arraystretch}{1.5}
\begin{tabular}{@{}P{2.8cm}P{5.8cm}@{}}
\toprule
\textbf{Alle Grade gerade} & Ein Eulerkreis existiert; er ist die Lösung, jede Straße genau einmal.\\
\textbf{$2k$ ungerade Ecken} & Man muss $k$ Wege doppelt gehen --- und zwar so, dass die Gesamtlänge der Doppelungen minimal wird.\\
\bottomrule
\end{tabular}
\end{center}

Die zweite Zeile führt auf eine Aufgabe, die wir erst in Kapitel~11 lösen können: Man muss die ungeraden Ecken paarweise so zuordnen, dass die Summe der Abstände kleinstmöglich ist --- eine \kursiv{Paarung} minimalen Gewichts.

\begin{center}
\begin{minipage}{0.88\textwidth}
\kursiv{Das Problem wurde 1962 von \Person{Mei-Ko Kwan} formuliert und heißt deshalb im Englischen \enquote{Chinese Postman Problem}. Es ist eines der wenigen Rundreiseprobleme, für die ein schnelles Lösungsverfahren bekannt ist --- der scheinbar ähnliche Fall des nächsten Kapitels ist es nicht.}
\end{minipage}
\end{center}
\end{merksatz}

\section{Übungsaufgaben}

\begin{aufgabe}
\boxtitel{Aufgabe 6.1 --- Das Haus vom Nikolaus}
Das bekannte Zeichenspiel verlangt, die abgebildete Figur in einem Zug zu zeichnen, ohne abzusetzen und ohne eine Linie doppelt zu ziehen. Zeigen Sie, dass es geht, und bestimmen Sie alle möglichen Startpunkte.

\begin{center}
\begin{tikzpicture}[scale=1.3]
\node[knotenb] (a) at (-0.9,-0.9) {$a$};
\node[knotenb] (b) at (0.9,-0.9) {$b$};
\node[knotenb] (c) at (0.9,0.5) {$c$};
\node[knotenb] (d) at (-0.9,0.5) {$d$};
\node[knotenb] (e) at (0,1.5) {$e$};
\draw[kante] (a)--(b);\draw[kante] (b)--(c);\draw[kante] (c)--(d);\draw[kante] (d)--(a);
\draw[kante] (a)--(c);\draw[kante] (b)--(d);\draw[kante] (d)--(e);\draw[kante] (e)--(c);
\end{tikzpicture}
\end{center}

\loesung
\schritt{Schritt 1 --- Grade bestimmen}
\begin{center}
\renewcommand{\arraystretch}{1.42}
\begin{tabular}{@{}P{1.2cm}P{3.4cm}P{1.4cm}P{1.6cm}@{}}
\toprule
\textbf{Ecke} & \textbf{Nachbarn} & $\deg$ & \textbf{Parität}\\
\midrule
$a$ & $b, d, c$ & $3$ & \textcolor{akzent}{ungerade}\\
$b$ & $a, c, d$ & $3$ & \textcolor{akzent}{ungerade}\\
$c$ & $b, d, a, e$ & $4$ & gerade\\
$d$ & $c, a, b, e$ & $4$ & gerade\\
$e$ & $d, c$ & $2$ & gerade\\
\bottomrule
\end{tabular}
\end{center}
Probe: $3+3+4+4+2 = 16 = 2 \cdot 8$. \checkmark

\schritt{Schritt 2 --- Das Kriterium anwenden}
Der Graph ist zusammenhängend und hat \kursiv{genau zwei} Ecken ungeraden Grades, nämlich $a$ und $b$. Nach Satz 6.5 gibt es einen offenen Eulerzug, und er muss in einer der beiden ungeraden Ecken beginnen und in der anderen enden.

\schritt{Schritt 3 --- Antwort auf die Frage nach den Startpunkten}
Die möglichen Startpunkte sind genau $\mathbf{a}$ und $\mathbf{b}$ --- die beiden unteren Ecken. Von $c$, $d$ oder $e$ aus ist es unmöglich.

\schritt{Schritt 4 --- Ein konkreter Zug}
Wir starten in $a$:
\[ a \to b \to c \to a \to d \to b \quad ? \]
Halt: Nach $b$ wären die Kanten $ab$, $bc$, $ca$, $ad$, $db$ verbraucht --- das sind fünf. An $b$ hängen $ab$, $bc$, $bd$, alle drei sind benutzt. Wir stecken fest, obwohl $de$, $ec$, $cd$ übrig sind. Der Zug war zulässig, aber ungeschickt.

Zweiter Versuch:
\[ a \to b \to c \to e \to d \to c \to a \to d \to b \]
Kanten der Reihe nach: $ab$, $bc$, $ce$, $ed$, $dc$, $ca$, $ad$, $db$. Das sind acht verschiedene --- alle. \checkmark\ Der Zug endet in $b$, wie vorhergesagt.

\begin{center}
\begin{minipage}{0.86\textwidth}
\kursiv{Der erste Versuch zeigt etwas Wichtiges: Satz 6.5 garantiert die \kursiv{Existenz}, nicht dass jedes Vorgehen zum Ziel führt. Wer sich verläuft, muss zurück. Das Verfahren von Hierholzer aus dem Beweis vermeidet das systematisch --- es baut Teilkreise und setzt sie ineinander, statt vorwärts zu raten.}
\end{minipage}
\end{center}
\end{aufgabe}

\begin{aufgabe}
\boxtitel{Aufgabe 6.2 --- Wann ist $K_n$ eulersch?}
Für welche $n$ besitzt der vollständige Graph $K_n$ einen Eulerkreis? Für welche einen offenen Eulerzug?

\loesung
\schritt{Schritt 1 --- Die Grade}
In $K_n$ ist jede Ecke mit allen $n-1$ übrigen benachbart, also $\deg(v) = n-1$ für alle $v$.

\schritt{Schritt 2 --- Eulerkreis}
Verlangt sind lauter gerade Grade, also muss $n-1$ gerade sein, also $n$ ungerade. Zusätzlich muss $K_n$ zusammenhängend sein und mindestens eine Kante haben --- das ist für $n \ge 2$ erfüllt.

\textbf{Antwort:} $K_n$ ist eulersch genau für ungerade $n \ge 3$, also $n = 3, 5, 7, 9, \ldots$

\schritt{Schritt 3 --- Offener Eulerzug}
Verlangt sind genau zwei ungerade Ecken. In $K_n$ haben aber \kursiv{alle} $n$ Ecken denselben Grad $n-1$. Sind alle ungerade, so sind es $n$ Stück, und $n = 2$ ist die einzige Möglichkeit.

Für $n = 2$: zwei Ecken vom Grad $1$, eine Kante. Der Zug ist $1 \to 2$. \checkmark

\textbf{Antwort:} Nur $K_2$.

\schritt{Probe an Beispielen}
\begin{center}
\renewcommand{\arraystretch}{1.42}
\begin{tabular}{@{}P{1.2cm}P{1.6cm}P{1.6cm}P{3.0cm}@{}}
\toprule
$n$ & $\deg$ & $m$ & \textbf{Befund}\\
\midrule
$3$ & $2$ & $3$ & Eulerkreis: das Dreieck\\
$4$ & $3$ & $6$ & vier ungerade Ecken --- nichts\\
$5$ & $4$ & $10$ & Eulerkreis mit $10$ Kanten\\
$6$ & $5$ & $15$ & sechs ungerade Ecken --- nichts\\
\bottomrule
\end{tabular}
\end{center}
\end{aufgabe}

\begin{aufgabe}
\boxtitel{Aufgabe 6.3 --- Ein Zerlegungssatz}
Zeigen Sie: Ein Graph, in dem jede Ecke geraden Grad hat, lässt sich in kantendisjunkte Kreise zerlegen --- das heißt, die Kantenmenge lässt sich so in Gruppen aufteilen, dass jede Gruppe ein Kreis ist.

\loesung
\schritt{Schritt 1 --- Induktion über die Kantenzahl}
\textbf{Anfang $m = 0$:} Die leere Zerlegung leistet das Verlangte.

\schritt{Schritt 2 --- Induktionsschritt}
Sei $m \ge 1$ und die Behauptung für alle Graphen mit weniger Kanten und lauter geraden Graden bewiesen.

\schritt{Schritt 3 --- Einen Kreis finden}
Da $m \ge 1$ gibt es eine Ecke mit Grad $\ge 1$, wegen der Parität also mit Grad $\ge 2$. Nach dem Hilfssatz aus dem Beweis von Satz 6.4 enthält $G$ einen Kreis $C$.

\schritt{Schritt 4 --- Ihn herausnehmen}
Sei $G' = G - E(C)$, also $G$ ohne die Kanten von $C$ (die Ecken bleiben).

Ein Kreis berührt jede Ecke entweder gar nicht oder genau zweimal --- er kommt an und geht weg. Also senkt das Entfernen von $C$ jeden Grad um $0$ oder $2$. Gerade minus gerade ist gerade: In $G'$ sind alle Grade weiterhin gerade.

\schritt{Schritt 5 --- Induktionsannahme}
$G'$ hat $m - |C| < m$ Kanten und lauter gerade Grade. Nach Induktionsannahme zerfällt $E(G')$ in kantendisjunkte Kreise.

\schritt{Schritt 6 --- Zusammensetzen}
Fügt man $C$ zu dieser Zerlegung hinzu, so erhält man eine Zerlegung von $E(G)$: Jede Kante liegt entweder in $C$ oder in $G'$, nie in beiden. $\blacksquare$

\schritt{Nachbemerkung --- der Zusammenhang zu Satz 6.4}
Der Satz gilt \kursiv{ohne} Zusammenhangsvoraussetzung. Für einen Eulerkreis braucht man zusätzlich, dass alles in einer Komponente liegt --- erst dann lassen sich die Kreise zu einem einzigen Zug verschmelzen. Das ist genau Schritt 4 des Hierholzer-Beweises.
\end{aufgabe}

\begin{geschichte}
\boxtitelg{Was Euler wirklich bewies --- und was Hierholzer nachtrug}
Eulers Aufsatz von 1736 enthält den notwendigen Teil vollständig: Er zeigt, dass ungerade Grade den Zug verhindern, und beantwortet damit die Königsberger Frage.

Den umgekehrten Teil --- dass gerade Grade genügen --- behauptete er, ohne ihn auszuführen. Er hielt ihn offenbar für offensichtlich. Er ist es nicht: Der Beweis braucht das Verschmelzungsargument aus Schritt 5, und das ist alles andere als selbstverständlich.

Den vollständigen Beweis lieferte \Person{Carl Hierholzer} erst 1873, hundertsiebenunddreißig Jahre später --- und auch das nur beinahe: Hierholzer starb 1871 im Alter von neunundzwanzig Jahren, bevor er die Arbeit niederschreiben konnte. Zwei Kollegen rekonstruierten sie aus seinen Vorträgen und veröffentlichten sie posthum.

\begin{center}
\begin{minipage}{0.9\textwidth}
\kursiv{Der Fall ist ein Musterbeispiel für die Warnung aus dem Vorwort: Ein Bild kann so überzeugend sein, dass selbst Euler den fehlenden Schritt nicht bemerkt. Wer heute schreibt \enquote{die Umkehrung ist klar}, steht in dieser Tradition --- und liegt womöglich ebenso daneben.}
\end{minipage}
\end{center}
\end{geschichte}

\begin{forschung}
\boxtitelp{Eulerzüge in der Sequenzierung des Erbguts}
Beim Zusammensetzen eines Genoms aus kurzen Leseabschnitten wird ein Graph gebaut, in dem die Ecken kurze Zeichenketten fester Länge sind und eine Kante zwei überlappende Ketten verbindet --- der \Begriff{De-Bruijn-Graph}. Die gesuchte Gesamtsequenz entspricht dann einem Eulerzug.

Dass man hier den Euler- und nicht den Hamiltonweg sucht, ist der entscheidende Kunstgriff. Der naheliegende Ansatz --- jede Lesung als Ecke, gesucht ein Weg durch alle Ecken --- führt auf ein Hamiltonproblem und ist praktisch unlösbar. Die Umstellung auf Kanten macht dieselbe Aufgabe schnell lösbar.

\begin{center}
{\small
\renewcommand{\arraystretch}{1.42}
\begin{tabular}{@{}P{3.0cm}P{2.2cm}P{2.8cm}@{}}
\toprule
\textbf{Ansatz} & \textbf{Gesucht} & \textbf{Aufwand}\\
\midrule
Überlappungs\-graph & Hamiltonweg & \textcolor{akzent}{praktisch unlösbar}\\
De-Bruijn-Graph & Eulerzug & \textcolor{gruen}{$\Oh(m)$}\\
\bottomrule
\end{tabular}}
\end{center}

\begin{center}
\begin{minipage}{0.88\textwidth}
\kursiv{Der Ansatz geht auf Pevzner, Tang und Waterman (2001) zurück und ist heute die Grundlage der meisten Zusammensetzungsverfahren. In der Praxis stören Lesefehler und Wiederholungen im Erbgut die Eindeutigkeit erheblich --- ein Genom hat typischerweise sehr viele Eulerzüge, und die Auswahl des richtigen ist der eigentlich schwierige Teil.}
\end{minipage}
\end{center}
\end{forschung}

\begin{merksatz}
\boxtitel{Was Sie jetzt können}
\begin{center}
\renewcommand{\arraystretch}{1.5}
\begin{tabular}{@{}P{3.0cm}P{5.6cm}@{}}
\toprule
\textbf{Multigraph} & Die Erweiterung von Definition 1.1, mit Schlingen und Parallelen\\
\textbf{Satz 6.4} & Eulerkreis $\iff$ zusammenhängend und alle Grade gerade\\
\textbf{Satz 6.5} & Offener Eulerzug $\iff$ genau zwei ungerade Ecken\\
\textbf{Hierholzer} & Ein Verfahren, das den Zug wirklich baut\\
\textbf{Königsberg} & Vier ungerade Ecken --- der Spaziergang existiert nicht\\
\bottomrule
\end{tabular}
\end{center}

\begin{center}
\begin{minipage}{0.88\textwidth}
\kursiv{Das nächste Kapitel stellt dieselbe Frage für Ecken statt Kanten --- und kommt zu einem völlig anderen Ergebnis.}
\end{minipage}
\end{center}
\end{merksatz}
